% ------------------------------------------------------------------------
% file `3_uaa1_fonctions-ineqations-systeme_02_2023-2024-appliquer-9-exercise-body.tex'
%   in folder `xsim/'
%
%     exercise of type `appliquer' with id `9'
%
% generated by the `appliquer' environment of the
%   `xsim' package v0.21 (2022/02/12)
% from source `3_uaa1_fonctions-ineqations-systeme_02_2023-2024' on 2023/10/20 on line 152
% ------------------------------------------------------------------------
    Détermine l'expression analytique des fonctions suivantes grâce aux informations données. N'oublie pas d'écrire l'expression complète de la fonction.
    \begin{tasks}
        \task Le graphe de la fonction est une droite qui passe par les points \((3,1)\) et \((-3,-5)\).
        \task La fonction admet \(-2\) comme racine et son graphe passe par le point \((2,10)\).
        \task Un tableau de valeur de la fonction est le suivant :
        \begin{center}
            \begin{tblr}{hlines, vline{1,2,5}, columns={mode=math, c, 1cm}}
                x    & -2 & 4 & 10 \\
                f(x) & 5  & 8 & 11 \\
            \end{tblr}
        \end{center}
    \end{tasks}
